% Generated draft: requires mathematical and source review.
\documentclass[11pt]{article}
\usepackage[margin=25mm]{geometry}
\usepackage{fontspec}
\setmainfont{Latin Modern Roman}
\usepackage{amsmath,amssymb,mathtools,unicode-math}
\setmathfont{Latin Modern Math}
\usepackage{enumitem,needspace,xurl,hyperref}
\setlength{\parindent}{0pt}
\setlength{\parskip}{5pt}
\setlength{\emergencystretch}{3em}
\newcommand{\N}{\mathbb{N}}
\newcommand{\Z}{\mathbb{Z}}
\newcommand{\Q}{\mathbb{Q}}
\newcommand{\R}{\mathbb{R}}
\newcommand{\C}{\mathbb{C}}
\newcommand{\F}{\mathbb{F}}
\newcommand{\A}{\mathbb{A}}
\newcommand{\PP}{\mathbb{P}}
\newcommand{\E}{\mathbb{E}}
\newcommand{\eps}{\varepsilon}
\newcommand{\dd}{\,\mathrm{d}}
\begin{document}
\section*{UnsolvedMath ID: 800059}
\textbf{Source problem code:} catalogue entry
\subsection{Definitions and mathematical statement}
The supplied question, attributed to Problem 8 of Hirzebruch's 1954 problem list, asks whether the Hodge numbers \(h^{p,q}\) and the Chern characteristic numbers of an algebraic variety \(V_n\) are topological invariants, and, if they are not, asks for all linear combinations of these numbers that are topological invariants.

The source excerpt does not specify the base field, smoothness, compactness, or the meaning of topological equivalence. The following makes the question precise conditionally in the setting of smooth projective complex varieties; this setting is not asserted to be the complete intended scope of the source.

For a smooth projective complex variety \(V\) of complex dimension \(n\), define
\[
h^{p,q}(V)=\dim_{\C}H^q(V,\Omega_V^p),\qquad 0\leq p,q\leq n,
\]
where \(\Omega_V^p\) is the sheaf of holomorphic differential \(p\)-forms. Write \(c_i(TV)\in H^{2i}(V;\Z)\) for the Chern classes of the holomorphic tangent bundle. For each tuple \(\alpha=(\alpha_1,\ldots,\alpha_n)\in\Z_{\geq0}^n\) satisfying \(\sum_{i=1}^n i\alpha_i=n\), define the Chern characteristic number
\[
C_\alpha(V)=\left\langle\prod_{i=1}^n c_i(TV)^{\alpha_i},[V]\right\rangle,
\]
where \([V]\) is the fundamental class for the complex orientation.

Fix a dimension \(n\), a coefficient field \(K\), and an equivalence relation \(\sim\) on the underlying real manifolds. These choices require clarification from the source. A numerical expression \(I\) is invariant under \(\sim\) if
\[
V\sim W\quad\Longrightarrow\quad I(V)=I(W)
\]
for every pair of varieties \(V,W\) in the specified class. The conditional classification problem is to determine all coefficient families \((a_{p,q},b_\alpha)\) in \(K\) such that
\[
L(V)=\sum_{p,q=0}^{n}a_{p,q}h^{p,q}(V)
+\sum_{\substack{\alpha\in\Z_{\geq0}^{n}\\\sum_{i=1}^{n}i\alpha_i=n}}b_\alpha C_\alpha(V)
\]
satisfies
\[
\forall V,W,\qquad V\sim W\ \Longrightarrow\ L(V)=L(W).
\]
The individual-invariance question asks which of the functions \(h^{p,q}\) and \(C_\alpha\) satisfy the same implication. The source does not settle whether \(\sim\) means homeomorphism, orientation-preserving homeomorphism, diffeomorphism, or orientation-preserving diffeomorphism; these give distinct formulations. It also does not specify whether mixed linear combinations of both kinds of numbers, or separate classifications for each kind, are intended. Singular or noncompact varieties require additional definitions before this formulation can be extended to them.
\subsection{Short English statement}
Determine which Hodge numbers, Chern characteristic numbers, and linear combinations of these numbers depend only on the topology of an algebraic variety. A precise version requires specifying the class of varieties, the coefficients, and whether topology includes an orientation. The displayed formulation is conditional on a smooth projective complex setting and requires human mathematical review.
\subsection{Sources}
\begin{itemize}
\item Supplied local UnsolvedMath registry record: permanent ID \(800059\), exact source problem code \texttt{catalogue entry}, title ``Hirzebruch's 1954 Problem List --- Problem 8,'' in the AMR Open Problem Lists collection. The registry lists Kotschick as proposer and \(2013\) as the proposed year; the title attributes the underlying problem to Hirzebruch's 1954 list.
\item Supplied source reference: Kotschick, \emph{Updates on Hirzebruch's 1954 Problem List} (2013), Problem 8, \url{https://arxiv.org/abs/1305.4623}. The supplied provenance reports access on \texttt{2026-10-16} and extraction method \texttt{arxiv\_latex}; neither access nor the source contents have been independently verified here.
\item The supplied registry labels the problem \texttt{open} and records an update on \texttt{2026-10-16}. This is supplied dated registry metadata, not an independently verified assessment of the mathematical literature. No separate literature-verification date or detailed literature assessment is supplied.
\end{itemize}
\end{document}
